\documentclass[12pt,a4paper]{jsarticle}
\usepackage[dvipdfm]{geometry}
\usepackage{xcolor}
\usepackage{graphicx}
\usepackage{array}
\usepackage{listings}
\usepackage{hyperref}
\usepackage{setspace}
\usepackage{multirow}
\usepackage{tcolorbox}

% ================== ページ設定 ==================
\geometry{top=30mm,bottom=30mm,left=18mm,right=18mm}
\pagestyle{plain}
\setstretch{1.2}

% ================== ユーザー定義コマンド ==================
% 日付欄の簡単入力: \pracdate{年}{月}{日}
\newcommand{\pracdate}[3]{#1年~#2月~#3日}
% または手動入力: \pracdate{2026}{4}{24}
\newcommand{\TODO}[1]{\textcolor{red}{【TODO: #1】}}
\newcommand{\NOTE}[1]{\textcolor{blue}{【注: #1】}}

% コード表示設定
\lstset{
    language=bash,
    basicstyle=\ttfamily\small,
    breaklines=true,
    frame=single,
    numbers=left,
    numberstyle=\tiny,
    keywordstyle=\color{blue},
    commentstyle=\color{gray},
    stringstyle=\color{red},
    backgroundcolor=\color{white},
    showstringspaces=false,
    tabsize=4
}

\begin{document}

% ==================== ページ1: 表紙 ====================
\thispagestyle{empty}

\vspace{1.5cm}

% ========== タイトル ==========
\centering
{\fontsize{32pt}{40pt}\selectfont\bfseries 情報科学実習I}\par
\vspace{0.5cm}
{\fontsize{32pt}{40pt}\selectfont\bfseries シェルスクリプト入門}\par
\vspace{0.5cm}
{\fontsize{32pt}{40pt}\selectfont\bfseries レポート}\par

\vspace{2.5cm}

% ========== 日付情報（カスタマイズ用） ==========
% 以下の行を編集して日付を入力してください
\raggedright
\setlength{\tabcolsep}{0pt}
\begin{tabular}{@{\hspace{1.5cm}}l@{\hspace{0.5cm}}l@{\hspace{2cm}}l}
実習実施日 & 1回目 & \pracdate{2026}{}{} \\[0.2cm]
& 2回目 & \pracdate{2026}{}{} \\[0.4cm]
レポート提出日 & & \pracdate{2026}{}{} \\[0.4cm]
再提出日 & & \pracdate{2026}{}{} \\
\end{tabular}

\vspace{2.5cm}

% ========== 提出者情報 ==========
\setlength{\tabcolsep}{0pt}
\begin{tabular}{@{\hspace{1.5cm}}l@{\hspace{0.5cm}}l}
{\textcolor{red}\bfseries レポート提出者} & \\[0.3cm]
学生番号 & \underline{\phantom{XXXXXXXXXXXXXXXXXXXX}} \\[0.2cm]
名前 & \underline{\phantom{XXXXXXXXXXXXXXXXXXXXXXXXXXXX}} \\
\end{tabular}

% ==================== ページ2以降: 本文 ====================
\newpage

\section*{1. 実習の目的}
本実習全体の目的を、自分の言葉で2～4文程度でまとめる。\par
\vspace{1.2cm}

\noindent
\underline{\hspace{15cm}}
\vspace{0.4cm}

\noindent
\underline{\hspace{15cm}}
\vspace{0.4cm}

\noindent
\underline{\hspace{15cm}}

\section*{2. 実習環境}
使用したOS、シェル、エディタ、参考資料、実行環境などを記入する。必要に応じて、授業スライドやテキストを参照したことも明記する。\par
\vspace{0.6cm}

\noindent
\begin{tabular}{|l|l|}
\hline
\textbf{項目} & \textbf{詳細} \\
\hline
OS & Ubuntu 24.04.3 LTS (Noble Numbat) \\
\hline
カーネル & Linux 6.14.0-37-generic \\
\hline
アーキテクチャ & x86\_64 \\
\hline
CPU & Intel(R) Core(TM) i5-14500T（論理CPU数：20） \\
\hline
シェル & tcsh 6.24.10 \\
\hline
bash バージョン & GNU bash 5.2.21 \\
\hline
ホスト名 & eib-pc064 \\
\hline
ロケール & ja\_JP.UTF-8 \\
\hline
タイムゾーン & Asia/Tokyo \\
\hline
テキスト/スライド & あり（必要に応じてファイル名を記入） \\
\hline
\end{tabular}

\vspace{0.6cm}

\noindent
補足:\quad 
\underline{\hspace{14cm}}

\vspace{1.5cm}

\section*{3. 実習方法}

本節には、各課題に共通する作業手順のみを簡潔に記載する（個別の操作内容・実行結果・考察は 4章に記載）。\par
\vspace{0.4cm}

\begin{enumerate}
    \item テキストおよび配布スライドで課題内容を確認した。
    \item エディタでシェルスクリプトを作成し、必要に応じて実行権限を付与した。
    \item 端末からスクリプトを実行し、引数を変えて動作を確認した。
    \item 実行結果を保存し、結果と考察を各課題欄（4章以降）に整理した。
\end{enumerate}

\noindent
※ 重複を避けるため、個別課題の詳細手順は 4章「結果と考察」に統合して記載する。

\newpage

\section*{4. 結果と考察}

各実習について、以下の順で記入する。
\begin{itemize}
    \item 目的
    \item スクリプト（必要ならコードを貼る）
    \item 実行例
    \item 結果の説明・考察
\end{itemize}

% ========== 実習1 ==========
\subsection*{4.1 実習1}

実習1は指定問題（問題2, 4, 5, 6, 7, 8, 10, 11）ごとに記入する。

\subsubsection*{問題2}
コード:\par
\begin{lstlisting}
#!/bin/sh
# ここに問題2のコード
\end{lstlisting}
出力例:\par
\begin{lstlisting}
# ここに実行コマンドと出力
\end{lstlisting}
方法・考察:\par
\underline{\hspace{15cm}}\\[0.4cm]
\underline{\hspace{15cm}}

\vspace{0.6cm}

\subsubsection*{問題4}
コード:\par
\begin{lstlisting}
#!/bin/sh
# ここに問題4のコード
\end{lstlisting}
出力例:\par
\begin{lstlisting}
# ここに実行コマンドと出力
\end{lstlisting}
方法・考察:\par
\underline{\hspace{15cm}}\\[0.4cm]
\underline{\hspace{15cm}}

\vspace{0.6cm}

\subsubsection*{問題5}
コード:\par
\begin{lstlisting}
#!/bin/sh
# ここに問題5のコード
\end{lstlisting}
出力例:\par
\begin{lstlisting}
# ここに実行コマンドと出力
\end{lstlisting}
方法・考察:\par
\underline{\hspace{15cm}}\\[0.4cm]
\underline{\hspace{15cm}}

\vspace{0.6cm}

\subsubsection*{問題6}
コード:\par
\begin{lstlisting}
#!/bin/sh
# ここに問題6のコード
\end{lstlisting}
出力例:\par
\begin{lstlisting}
# ここに実行コマンドと出力
\end{lstlisting}
方法・考察:\par
\underline{\hspace{15cm}}\\[0.4cm]
\underline{\hspace{15cm}}

\vspace{0.6cm}

\subsubsection*{問題7}
コード:\par
\begin{lstlisting}
#!/bin/sh
# ここに問題7のコード
\end{lstlisting}
出力例:\par
\begin{lstlisting}
# ここに実行コマンドと出力
\end{lstlisting}
方法・考察:\par
\underline{\hspace{15cm}}\\[0.4cm]
\underline{\hspace{15cm}}

\vspace{0.6cm}

\subsubsection*{問題8}
コード:\par
\begin{lstlisting}
#!/bin/sh
# ここに問題8のコード
\end{lstlisting}
出力例:\par
\begin{lstlisting}
# ここに実行コマンドと出力
\end{lstlisting}
方法・考察:\par
\underline{\hspace{15cm}}\\[0.4cm]
\underline{\hspace{15cm}}

\vspace{0.6cm}

\subsubsection*{問題10}
コード:\par
\begin{lstlisting}
#!/bin/sh
# ここに問題10のコード
\end{lstlisting}
出力例:\par
\begin{lstlisting}
# ここに実行コマンドと出力
\end{lstlisting}
方法・考察:\par
\underline{\hspace{15cm}}\\[0.4cm]
\underline{\hspace{15cm}}

\vspace{0.6cm}

\subsubsection*{問題11}
コード:\par
\begin{lstlisting}
#!/bin/sh
# ここに問題11のコード
\end{lstlisting}
出力例:\par
\begin{lstlisting}
# ここに実行コマンドと出力
\end{lstlisting}
方法・考察:\par
\underline{\hspace{15cm}}\\[0.4cm]
\underline{\hspace{15cm}}

\vspace{1.0cm}

% ========== 実習2 ==========
\subsection*{4.2 実習2}

\subsubsection*{目的}
\underline{\hspace{15cm}}

\vspace{0.8cm}

\noindent
\begin{minipage}[t]{\linewidth}
\vspace{0.2cm}
\underline{\hspace{15cm}}\\[0.6cm]
\underline{\hspace{15cm}}\\[0.6cm]
\underline{\hspace{15cm}}\\[0.6cm]
\underline{\hspace{15cm}}\\[0.6cm]
\underline{\hspace{15cm}}
\end{minipage}

\subsubsection*{スクリプト}
\begin{lstlisting}
#!/bin/bash
# スクリプトをここに記入

\end{lstlisting}

\subsubsection*{実行例と説明}
実行例:\par
\begin{lstlisting}
# ここに実行コマンドと出力を記入
# 例: ./script.sh arg1 arg2
#     出力...
\end{lstlisting}
結果・考察:\par
\underline{\hspace{15cm}}\\[0.4cm]
\underline{\hspace{15cm}}\\[0.4cm]
\underline{\hspace{15cm}}\par

\vspace{1.5cm}

% ========== 実習3 ==========
\subsection*{4.3 実習3}

\subsubsection*{目的}
\underline{\hspace{15cm}}

\vspace{0.8cm}

\subsubsection*{スクリプト}
\begin{lstlisting}
#!/bin/bash
# スクリプトをここに記入

\end{lstlisting}

\subsubsection*{実行例と説明}
実行例:\par
\begin{lstlisting}
# ここに実行コマンドと出力を記入
# 例: ./script.sh arg1 arg2
#     出力...
\end{lstlisting}
結果・考察:\par
\underline{\hspace{15cm}}\\[0.4cm]
\underline{\hspace{15cm}}\\[0.4cm]
\underline{\hspace{15cm}}\par

\vspace{1.5cm}

% ========== 実習4 ==========
\subsection*{4.4 実習4}

実習4は「同じコードを引数だけ変えて実行する」形式なので、コードは1つ、実行例を複数記入する。

\subsubsection*{対象問題（例：実習4の該当番号）}
問題番号:\quad \underline{\hspace{5cm}}\par

コード:\par
\begin{lstlisting}
#!/bin/sh
# ここに実習4のコード（1つ）
\end{lstlisting}

実行例1（引数パターン1）:\par
\begin{lstlisting}
# 例: ./j4.sh 13
# 出力:
\end{lstlisting}

実行例2（引数パターン2）:\par
\begin{lstlisting}
# 例: ./j4.sh 10
# 出力:
\end{lstlisting}

実行例3（引数パターン3）:\par
\begin{lstlisting}
# 例: ./j4.sh 2
# 出力:
\end{lstlisting}

方法・考察:\par
\underline{\hspace{15cm}}\\[0.4cm]
\underline{\hspace{15cm}}\\[0.4cm]
\underline{\hspace{15cm}}

\newpage

\section*{5. 実習操作その2}

実習5、および実習6～9のうち選んだ題材について、同様に「目的」「方法」「スクリプト」「実行例」「考察」を記入する。\par
\vspace{0.5cm}

\subsection*{5.1 実習5}
実習5（外部コマンド・リダイレクト・sed/sort など）は複数問題があるため、問題番号ごとに記入する。

\subsubsection*{問題 19}
目的:\par
\underline{\hspace{15cm}}

方法:\par
\underline{\hspace{15cm}}

スクリプト/コマンド:\par
\begin{lstlisting}
# ここに問題19で試したコマンド
\end{lstlisting}

実行例:\par
\begin{lstlisting}
# ここに実行コマンドと出力
\end{lstlisting}

考察:\par
\underline{\hspace{15cm}}\\[0.4cm]
\underline{\hspace{15cm}}\\[0.4cm]
\underline{\hspace{15cm}}

\vspace{0.7cm}

\subsubsection*{発展 20（任意）}
目的:\par
\underline{\hspace{15cm}}

方法:\par
\underline{\hspace{15cm}}

スクリプト/コマンド:\par
\begin{lstlisting}
# ここにヒアドキュメント等の内容
\end{lstlisting}

実行例:\par
\begin{lstlisting}
# ここに実行コマンドと出力
\end{lstlisting}

考察:\par
\underline{\hspace{15cm}}\\[0.4cm]
\underline{\hspace{15cm}}\\[0.4cm]
\underline{\hspace{15cm}}

\vspace{0.7cm}

\subsubsection*{問題 21}
目的:\par
\underline{\hspace{15cm}}

方法:\par
\underline{\hspace{15cm}}

スクリプト/コマンド:\par
\begin{lstlisting}
# ここに問題21で試したコマンド
\end{lstlisting}

実行例:\par
\begin{lstlisting}
# ここに実行コマンドと出力
\end{lstlisting}

考察:\par
\underline{\hspace{15cm}}\\[0.4cm]
\underline{\hspace{15cm}}\\[0.4cm]
\underline{\hspace{15cm}}

\vspace{0.7cm}

\subsubsection*{問題 22}
目的:\par
\underline{\hspace{15cm}}

方法:\par
\underline{\hspace{15cm}}

スクリプト/コマンド:\par
\begin{lstlisting}
# ここに問題22で試したコマンド
\end{lstlisting}

実行例:\par
\begin{lstlisting}
# ここに実行コマンドと出力
\end{lstlisting}

考察:\par
\underline{\hspace{15cm}}\\[0.4cm]
\underline{\hspace{15cm}}\\[0.4cm]
\underline{\hspace{15cm}}

\vspace{0.7cm}

\subsubsection*{問題 23}
目的:\par
\underline{\hspace{15cm}}

方法:\par
\underline{\hspace{15cm}}

スクリプト/コマンド:\par
\begin{lstlisting}
#!/bin/sh
# ここに問題23のコード
\end{lstlisting}

実行例:\par
\begin{lstlisting}
# ここに実行コマンドと出力
\end{lstlisting}

考察:\par
\underline{\hspace{15cm}}\\[0.4cm]
\underline{\hspace{15cm}}\\[0.4cm]
\underline{\hspace{15cm}}

\vspace{1cm}

\subsection*{5.2 実習6～9（選択分）}
実習6～9は、実施した問題ごとにこのテンプレート一式をそのまま複製して記入する。
選んだ実習番号:\quad \underline{\hspace{6cm}}\par
目的:\quad \underline{\hspace{13cm}}\par
方法:\quad \underline{\hspace{13cm}}\par
スクリプト:\par
\begin{lstlisting}
#!/bin/bash
# ここに記入
\end{lstlisting}
実行例:\par
\begin{lstlisting}
# ここに実行コマンドと出力
\end{lstlisting}
考察:\par
\underline{\hspace{15cm}}\\[0.4cm]
\underline{\hspace{15cm}}\\[0.4cm]
\underline{\hspace{15cm}}\par

\newpage

% ==================== 参考文献 ====================
\section*{参考文献}
参考にしたテキスト、スライド、Webページなどを記載する。

\begin{thebibliography}{99}
\bibitem{ref1} 

\bibitem{ref2} 

\end{thebibliography}

\section*{備考}
本レポートの内容は自分で考えて作成した。LaTeXテンプレートの作成には GitHub Copilot for VS Code を使用した。

\end{document}
